\documentclass[a4paper,12pt]{article}
\usepackage[utf8]{inputenc}
\usepackage[T1]{fontenc}
\usepackage[ngerman]{babel}
\usepackage{geometry}
\usepackage{amsmath}
\usepackage{booktabs}
\usepackage{hyperref}
\geometry{margin=2.5cm}

\begin{document}

\section*{Zusammenfassung: Ashwin et al. (2024)}
\textbf{Titel:} Physics Informed Deep Learning for Flow and Force Predictions
in Dense Ellipsoidal Particle Suspensions \\
\textbf{Autoren:} Neil Raj Ashwin, Danesh Tafti, Nikhil Muralidhar, Ze Cao \\
\textbf{Journal:} Powder Technology 439 (2024) 119684

% ----------------------------------------------------------------
\subsection*{Problemstellung und Motivation}
% ----------------------------------------------------------------

Dichte Feststoff-Fluid-Suspensionen treten in zahlreichen industriellen und
natürlichen Prozessen auf, darunter Sedimentation, Wirbelschichtreaktoren und
chemische Verarbeitungsanlagen. Das zentrale numerische Hilfsmittel zur genauen
Untersuchung solcher Systeme sind \emph{Particle Resolved Simulations} (PRS), die
die Strömung um jedes einzelne Partikel explizit auflösen. Aufgrund des hohen
Rechenaufwands ist PRS jedoch auf kleine Systeme mit typischerweise weniger als
$10^5$ Partikeln beschränkt. Gröbere Modelle wie Euler-Lagrange- oder
Zwei-Fluid-Modelle müssen die Kraft- und Drehmomentübertragung zwischen
Partikeln und Fluid durch phänomenologische Korrelationen annähern, die nur
die mittleren Bedingungen (Reynolds-Zahl, Feststoffanteil) berücksichtigen und
daher die individuellen Variationen der Widerstandskräfte zwischen Partikeln
nicht erfassen.

Die Autoren begegnen dieser Lücke, indem sie Deep-Learning-Surrogatmodelle
entwickeln, die auf PRS-Daten trainiert werden und in der Lage sind, sowohl die
interstitiellen Strömungsfelder als auch die Widerstandskräfte auf einzelne
Partikel in Echtzeit vorherzusagen.

% ----------------------------------------------------------------
\subsection*{Datengenerierung und Versuchsaufbau}
% ----------------------------------------------------------------

Als Grundlage dienen PRS-Simulationen ellipsoidaler Partikel mit einem
Achsenverhältnis von 2{,}5, die mit dem hauseigenen Code
\textit{GenIDLEST} unter Verwendung der \emph{Immersed Boundary Method} (IBM)
durchgeführt werden. Die dimensionslosen inkompressiblen Navier-Stokes-Gleichungen
werden gelöst:
\begin{equation}
    \frac{\partial u_i}{\partial x_i} = 0, \qquad
    \rho\frac{\partial u_i}{\partial t} + \rho\frac{\partial}{\partial x_j}(u_i u_j)
    = -\frac{\partial p}{\partial x_i}
      + \frac{\partial}{\partial x_j}\!\left[\frac{1}{Re}\frac{\partial u_i}{\partial x_j}\right].
\end{equation}
Der Datensatz umfasst drei Feststoffanteile ($\varphi = 0{,}1;\;0{,}2;\;0{,}3$)
und vier Reynolds-Zahlen ($Re = 10;\;50;\;100;\;200$), was insgesamt
4632 Partikel-Datenpunkte ergibt. Für jedes betrachtete Partikel
(\emph{particle of interest}, POI) wird ein kubischer Ausschnitt der Größe
$\pm2{,}5$ Einheiten in alle Raumrichtungen um das POI extrahiert und auf ein
Gitter von $51^3$ Voxeln subsampelt.

% ----------------------------------------------------------------
\subsection*{Methodik: U-Net-Architektur und Eingaberepräsentation}
% ----------------------------------------------------------------

\paragraph{Eingaberepräsentationen.}
Es werden zwei Eingabedarstellungen verglichen. Die \emph{binäre Methode}
weist jedem Voxel den Wert 0 (Feststoff) oder 1 (Fluid) zu. Die
\emph{5-Kanal-Distanzfunktion} kodiert für jeden Fluidvoxel die Abstände zur
nächstgelegenen, zweitnächstgelegenen, \ldots, fünftnächstgelegenen
Partikeloberfläche:
\begin{equation}
    D^n = \begin{cases}
        |\delta_n| & \text{Fluidvoxel} \\
        0          & \text{Feststoffvoxel}
    \end{cases}
\end{equation}
Der Vorteil dieses Ansatzes liegt darin, dass höhere Kanalindizes
Informationen über weiter entfernte Partikel tragen, die außerhalb des
$5^3$-Ausschnitts liegen und somit nicht direkt im Eingabetensor sichtbar sind.

\paragraph{U-Net-Architektur.}
Das U-Net besteht aus einem Encoder-Zweig (sukzessive 3D-Faltungen mit
Max-Pooling) und einem Decoder-Zweig (Upsampling mit bilinearer Interpolation).
Skip-Connections verbinden korrespondierende Encoder- und Decoder-Ebenen, um
hochauflösende Merkmale wiederzuverwenden und dem Vanishing-Gradient-Problem
entgegenzuwirken. Im Flaschenhals werden Reynolds-Zahl und Feststoffanteil als
skalare Eingaben in den Tensor eingebettet. Vier separate Decoder-Zweige
sagen gleichzeitig die $x$-, $y$-, $z$-Geschwindigkeiten und das Druckfeld
vorher. Als Verlustfunktion wird der mittlere quadratische Fehler (MSE) mit
dem Adam-Optimierer minimiert.

\paragraph{Druckvorverarbeitung.}
Eine wichtige Vorverarbeitungsentscheidung betrifft das Druckfeld: Da der
mittlere Druck in Strömungsrichtung stark abfällt und dieses globale Signal
die lokalen Druckstörungen der Partikel überdeckt, trainieren die Autoren auf
dem Fluktuationsdruckfeld
\begin{equation}
    p'(x,y,z) = p(x,y,z) - \bar{p}(x),
\end{equation}
wobei $\bar{p}(x)$ der ebenengemittelte Druck in der jeweiligen $x$-Ebene ist.
Dieses Vorgehen verbessert die Druckvorhersagen erheblich.

% ----------------------------------------------------------------
\subsection*{Ergebnisse: Strömungsfeldvorhersage}
% ----------------------------------------------------------------

Die 5-Kanal-Distanzfunktion übertrifft die binäre Eingabe konsistent für alle
drei Geschwindigkeitskomponenten. Der mittlere $R^2$ erreicht bei
$\varphi = 0{,}3$ und $Re = 10$ Werte von $\sim\!0{,}94$ für die
$x$-Geschwindigkeit und $\sim\!0{,}98$ für den Druck. Die Vorhersagegenauigkeit
sinkt mit steigender Reynolds-Zahl (zunehmende Nichtlinearitäten) und steigt
mit zunehmendem Feststoffanteil (dichtere Partikelanordnungen sind leichter zu
lernen). Interessanterweise zeigt sich, dass das Modell auf ungesehenen
Reynolds-Zahlen ($Re = 50$ und $100$) ähnlich gut abschneidet wie auf
Reynolds-Zahlen aus dem Trainingsdatensatz -- die Interpolation zwischen
Strömungsbedingungen ist leichter als die Generalisierung auf neue
Partikelanordnungen.

Die physikalische Konsistenz der vorhergesagten Felder wird anhand der
Kontinuitätsgleichung und der Impulsgleichung überprüft. Die Residuen sind
zentriert um Null und zeigen eine ähnliche Verteilung wie die Residuen der
(ebenfalls subgesampelten) PRS-Felder.

\paragraph{Widerstandskraftvorhersage.}
Die Autoren vergleichen zwei Strategien: (1) ein CNN, das direkt aus der
5-Kanal-Distanzfunktion die Widerstandskraft vorhersagt, und (2) ein CNN, das
die vorhergesagten Strömungsfelder des U-Nets als Eingabe verwendet. Methode~2
erzielt durchgehend bessere Ergebnisse ($R^2 = 0{,}927$ vs.\
$0{,}865$; $\sim\!85\%$ aller Partikel innerhalb eines relativen Fehlers von
$20\%$ vs.\ $71\%$ für Methode~1). Der Vorteil von Methode~2 ist besonders
ausgeprägt bei niedrigem Feststoffanteil ($\varphi = 0{,}1$), wo die
Datendichte gering ist -- das Modell profitiert in diesem Fall stark von den
zusätzlichen physikalischen Informationen der Strömungsfelder.

% ----------------------------------------------------------------
\subsection*{Relevanz für die Masterarbeit}
% ----------------------------------------------------------------

Dieses Paper ist in mehrfacher Hinsicht eng mit der Masterarbeit verknüpft:

\paragraph{Gemeinsamer Ansatz: U-Net zur Strömungsfeldvorhersage.}
Wie in der vorliegenden Arbeit wird eine U-Net-Architektur eingesetzt, um
vollständige Strömungsfelder aus geometrischen Eingabeinformationen
vorherzusagen. Ashwin et al.\ operieren in 3D mit der Navier-Stokes-Gleichung
bei moderaten Reynolds-Zahlen, während die Masterarbeit 2D-Stokes-Strömungen
($Re \to 0$) im Kontext magmatischer Sedimentation untersucht. In beiden Fällen
sind die langen hydrodynamischen Wechselwirkungen zwischen Partikeln das
zentrale Modellierungsproblem.

\paragraph{Eingaberepräsentation und physikalische Informationen.}
Die 5-Kanal-Distanzfunktion von Ashwin et al.\ stellt Partikelinformationen
jenseits des lokalen Ausschnitts zur Verfügung -- konzeptuell verwandt mit dem
binären Belegungsfeld, das in der Masterarbeit als Eingabe genutzt wird. Die
Überlegenheit der distanzbasierten Darstellung gegenüber der binären Methode
unterstützt die Entscheidung, physikalisch informierte Eingabemerkmale zu
verwenden.

\paragraph{Druckvorverarbeitung als Analogie zur Strom\-funktion.}
Die Subtraktion des mittleren Druckgradienten ($p \to p'$) hat ein direktes
Analogon in der Masterarbeit: Die Verwendung der Stromfunktion $\psi$ anstelle
der Geschwindigkeitskomponenten eliminiert die Divergenzbedingung
$\nabla \cdot \mathbf{u} = 0$ analytisch und reduziert damit die Freiheitsgrade
des Lernproblems -- dies entspricht dem Entfernen eines dominanten, schwer zu
lernenden Signals, das die eigentlich interessanten Partikelinteraktionen
überdecken würde.

\paragraph{Generalisierung auf ungesehene Konfigurationen.}
Der zentrale Befund von Ashwin et al., dass die Generalisierung auf neue
Partikelanordnungen deutlich schwieriger ist als die Interpolation auf neue
Strömungsbedingungen, hat direkte Implikationen für die Masterarbeit. Die
Kernfragestellung -- ob Modelle, die auf $N$ Kristalle trainiert wurden, auf
$N + m$ Kristalle generalisieren -- entspricht genau dem schwierigeren dieser
beiden Generalisierungsszenarien. Die Ergebnisse von Ashwin et al.\
unterstreichen, dass für diese Art der Generalisierung physikalisch reichhaltige
Eingaben und Ausgaberepräsentationen besonders wichtig sind.

\paragraph{Verwendung von Hilfsverlusten und mehrstufigem Training.}
Ashwin et al.\ zeigen, dass das Training auf Strömungsfeldern (Hilfsaufgabe)
und anschließende Kraft\-vorhersage eine erhebliche Verbesserung gegenüber dem
direkten Lernen der Zielvariable bringt -- insbesondere wenn die Datenlage dünn
ist. Dieser Ansatz des indirekten, physikalisch informierten Lernens ist
konzeptuell verwandt mit dem in der Masterarbeit untersuchten Residuallern\-ansatz,
bei dem $\Delta\mathbf{v}_\text{learned}$ als Korrektur zur analytischen
Stokes-Lösung gelernt wird.

\paragraph{Skalierbarkeit und Datenmenge.}
Mit 4632 Datenpunkten und 12 Experimenten operieren Ashwin et al.\ in einem
ähnlichen Datenregime wie die Masterarbeit (200--300 LaMEM-Simulationen).
Die vergleichbaren Trainingsgrößen bestätigen, dass physikalisch strukturierte
Modelle auch in datenknappen Szenarien erfolgreich sein können.

\end{document}